% AUTHORED
\section{Scaled dot-product self-attention}\label{sec:08}
\conceptbadge{08}

Self-attention lets each token decide which other tokens are relevant to its
next representation. Starting from token vectors $X$, learned linear projections
produce queries, keys, and values:
\[
  Q = XW_Q,\qquad K = XW_K,\qquad V = XW_V.
\]
A query asks what a token is looking for. A key describes what a token offers for
matching. A value is the information to be mixed if the match is strong.

\subsection*{Scores and weights}
For a sequence of length $T$ and key/query width $d_k$, the raw attention scores
are
\[
  S = \frac{QK^\top}{\sqrt{d_k}},
\]
with shape $T \times T$ for one sequence. Row $i$ contains how strongly token
$i$ attends to every token $j$. Applying softmax row by row gives attention
weights:
\[
  A_{ij} = \frac{\exp(S_{ij})}{\sum_{\ell=1}^{T}\exp(S_{i\ell})}.
\]
Each row of $A$ sums to one, so each output token receives a convex weighted
summary of value vectors:
\[
  O = AV.
\]

\subsection*{Why the scale is there}
If query and key components have roughly unit variance, their dot product grows
in variance with $d_k$. Large dot products make softmax very sharp, which can
produce tiny gradients early in training. Dividing by $\sqrt{d_k}$ keeps scores
in a more stable range.

\subsection*{Self-attention versus recurrence}
Self-attention compares all token positions in parallel. It does not need to
process tokens one at a time as an RNN would. The cost is the $T \times T$ score
matrix, which becomes expensive for long context lengths. In this course the
contexts are small, so clarity matters more than optimization.

\subsection*{What is still missing}
Plain self-attention allows every token to read every other token, including
future positions. That is fine for some encoders but wrong for autoregressive
language modeling. Concept 09 adds the causal mask so position $t$ can only use
positions $\le t$.

\subsection*{Exercise connection}
The exercise implements the unmasked formula
$\operatorname{softmax}(QK^\top/\sqrt{d_k})V$. P4 then integrates the same
operation into the cumulative language model.
